\documentclass[12pt, a4paper]{article}
\usepackage{../../../myconfig}

\begin{document}

\titlebox{Chủ đề: Oxyz}{Vecto và các phép toán vecto \\ trong không gian}{Oxyz}

% --- CÂU 1 ---
\questionLinesManual{5}{1}{Cho tứ diện \(ABCD\). Hỏi có bao nhiêu vectơ khác vectơ \(\vec{0}\) mà mỗi vectơ có điểm đầu, điểm cuối là hai đỉnh của tứ diện \(ABCD\)?}{%
    \begin{tasks}(4)
        \task 12
        \task 4
        \task 10
        \task 8
    \end{tasks}
}

% --- CÂU 2 ---
\questionLinesManual{5}{2}{Cho hình lập phương \(ABCD.A'B'C'D'\). Mệnh đề nào sau đây \textbf{sai}?}{%
    \begin{tasks}(2)
        \task \(\vec{AB} + \vec{AD} + \vec{AA'} = \vec{AC'}\)
        \task \(\vec{AC} = \vec{AB} + \vec{AD}\)
        \task \(\left| \vec{AB} \right| = \left| \vec{CD} \right|\)
        \task \(\vec{AB} = \vec{CD}\)
    \end{tasks}
}

% --- CÂU 3 ---
\questionLinesManual{5}{3}{Cho hình tứ diện \(ABCD\) có trọng tâm \(G\). Mệnh đề nào sau đây \textbf{sai}?}{%
    \begin{tasks}(2)
        \task \(\vec{GA} + \vec{GB} + \vec{GC} + \vec{GD} = \vec{0}\)
        \task \(\vec{OG} = \dfrac{1}{4}\left( \vec{OA} + \vec{OB} + \vec{OC} + \vec{OD} \right)\)
        \task \(\vec{AG} = \dfrac{2}{3}\left( \vec{AB} + \vec{AC} + \vec{AD} \right)\)
        \task \(\vec{AG} = \dfrac{1}{4}\left( \vec{AB} + \vec{AC} + \vec{AD} \right)\)
    \end{tasks}
}

% --- CÂU 4 ---
\questionLinesManual{5}{4}{Cho tứ diện \(ABCD\), gọi \(I, J\) lần lượt là trung điểm của \(AB\) và \(CD\); Đẳng thức nào \textbf{sai}?}{%
    \begin{tasks}(2)
        \task \(\vec{IJ} = \dfrac{1}{2}\left( \vec{AC} + \vec{BD} \right)\)
        \task \(\vec{IJ} = \dfrac{1}{2}\left( \vec{AD} + \vec{BC} \right)\)
        \task \(\vec{IJ} = \dfrac{1}{2}\left( \vec{DC} + \vec{AD} + \vec{BD} \right)\)
        \task \(\vec{IJ} = \dfrac{1}{2}\left( \vec{AB} + \vec{CD} \right)\)
    \end{tasks}
}

% --- CÂU 5 ---
\questionLinesManual{5}{5}{Cho hình hộp chữ nhật \(ABCD.A'B'C'D'\). Khi đó, vectơ bằng vectơ \(\vec{AB}\) là vectơ nào dưới đây?}{%
    \begin{tasks}(4)
        \task \(\vec{D'C'}\)
        \task \(\vec{BA}\)
        \task \(\vec{CD}\)
        \task \(\vec{B'A'}\)
    \end{tasks}
}

% --- CÂU 6 ---
\questionLinesManual{5}{6}{Cho hình chóp \(S.ABCD\) có đáy \(ABCD\) là hình bình hành. Đặt \(\vec{SA} = \vec{a}\); \(\vec{SB} = \vec{b}\); \(\vec{SC} = \vec{c}\); \(\vec{SD} = \vec{d}\). Khẳng định nào sau đây đúng?}{%
    \begin{tasks}(4)
        \task \(\vec{a} + \vec{b} + \vec{c} + \vec{d} = \vec{0}\)
        \task \(\vec{a} + \vec{b} = \vec{c} + \vec{d}\)
        \task \(\vec{a} + \vec{d} = \vec{b} + \vec{c}\)
        \task \(\vec{a} + \vec{c} = \vec{d} + \vec{b}\)
    \end{tasks}
}

% --- CÂU 7 ---
\questionLinesManual{5}{7}{Cho tứ diện \(ABCD\). Gọi \(P, Q\) là trung điểm của \(AB\) và \(CD\). Chọn khẳng định đúng?}{%
    \begin{tasks}(2)
        \task \(\vec{PQ} = \dfrac{1}{2}\left( \vec{BC} + \vec{AD} \right)\)
        \task \(\vec{PQ} = \dfrac{1}{2}\left( \vec{BC} - \vec{AD} \right)\)
        \task \(\vec{PQ} = \vec{BC} + \vec{AD}\)
        \task \(\vec{PQ} = \dfrac{1}{4}\left( \vec{BC} + \vec{AD} \right)\)
    \end{tasks}
}

% --- CÂU 8 ---
\questionLinesManual{5}{8}{Cho hình lăng trụ \(ABC.A'B'C'\). Đặt \(\vec{AB} = \vec{a}, \vec{AA'} = \vec{b}, \vec{AC} = \vec{c}\). Khẳng định nào sau đây đúng?
\begin{center}
    \begin{tikzpicture}[scale=1.2]
    % Định nghĩa tọa độ các đỉnh đáy dưới
    \coordinate (A) at (0, 0);
    \coordinate (B) at (0.8, -1.2);
    \coordinate (C) at (2.8, 0);

    % Vectơ tịnh tiến lên đáy trên (độ nghiêng và chiều cao)
    \coordinate (shift) at (0.2, 3.2);

    % Định nghĩa tọa độ các đỉnh đáy trên
    \coordinate (A') at ($(A) + (shift)$);
    \coordinate (B') at ($(B) + (shift)$);
    \coordinate (C') at ($(C) + (shift)$);

    % Vẽ các cạnh khuất (nét đứt)
    \draw[thick, dashed, blue!70!black] (A) -- (C);

    % Vẽ các cạnh nhìn thấy của đáy dưới (nét liền)
    \draw[very thick, blue!70!black] (A) -- (B) -- (C);

    % Vẽ các cạnh bên (nét liền)
    \draw[very thick, blue!70!black] (A) -- (A');
    \draw[very thick, blue!70!black] (B) -- (B');
    \draw[very thick, blue!70!black] (C) -- (C');

    % Vẽ các cạnh của đáy trên (nét liền)
    \draw[very thick, blue!70!black] (A') -- (B') -- (C') -- cycle;

    % Vẽ các điểm (nút) màu đỏ tại các đỉnh
    \foreach \p in {A, B, C, A', B', C'} {
        \fill[red] (\p) circle (2pt);
    }

    % Đặt nhãn cho các đỉnh
    \node[left] at (A) {$A$};
    \node[below] at (B) {$B$};
    \node[right] at (C) {$C$};
    \node[above left] at (A') {$A'$};
    \node[below left] at (B') {$B'$};
    \node[above right] at (C') {$C'$};

\end{tikzpicture}
\end{center}
}{%
    \begin{tasks}(2)
        \task \(\vec{B'C} = \vec{a} + \vec{b} - \vec{c}\)
        \task \(\vec{B'C} = -\vec{a} + \vec{b} - \vec{c}\)
        \task \(\vec{B'C} = -\vec{a} - \vec{b} + \vec{c}\)
        \task \(\vec{B'C} = -\vec{a} + \vec{b} + \vec{c}\)
    \end{tasks}
}
% --- CÂU 9 ---
\questionLinesManual{5}{9}{Cho hình lập phương \(ABCD.A'B'C'D'\). Tính \(\cos\left( \vec{BD}, \vec{A'C'} \right)\)}{%
    \begin{tasks}(4)
        \task \(\cos\left( \vec{BD}, \vec{A'C'} \right) = 0\)
        \task \(\cos\left( \vec{BD}, \vec{A'C'} \right) = 1\)
        \task \(\cos\left( \vec{BD}, \vec{A'C'} \right) = \dfrac{1}{2}\)
        \task \(\cos\left( \vec{BD}, \vec{A'C'} \right) = \dfrac{\sqrt{2}}{2}\)
    \end{tasks}
}

% --- CÂU 10 ---
\questionLinesManual{5}{10}{Cho tứ diện \(ABCD\) và điểm \(G\) thỏa \(\vec{GA} + \vec{GB} + \vec{GC} = \vec{0}\) (\(G\) là trọng tâm của tứ diện). Gọi \(O\) là giao điểm của \(GA\) và mặt phẳng \((BCD)\). Trong các khẳng định sau, khẳng định nào sai?}{%
    \begin{tasks}(4)
        \task \(\vec{GA} = -2\vec{OG}\)
        \task \(\vec{GA} = 4\vec{OG}\)
        \task \(\vec{GA} = 3\vec{OG}\)
        \task \(\vec{GA} = 2\vec{OG}\)
    \end{tasks}
}


\end{document}